\section{Finite-data common-map identification}
\label{sec:framework}

\subsection{Transition channels and admissible maps}

Let $\A$ be a finite-dimensional $C^*$-algebra representing the retained coarse record, equipped with its physical trace $\tau$.  A complete source experiment at time $t$ is represented by a channel
\begin{equation}
\Lambda_t:M_2\longrightarrow \A,
\end{equation}
where the two-dimensional source is probed on all four matrix units, including its off-diagonal coherences.  For a sampled transition $j$ from $t_j$ to $t_j+\Delta$, write
\begin{equation}
\Lambda_j^{\mathrm{cur}}=\Lambda_{t_j},\qquad
\Lambda_j^{\mathrm{next}}=\Lambda_{t_j+\Delta}.
\end{equation}

The admissible class $\C$ consists of physical CPTP maps
\begin{equation}
\Phi:\A\to\A
\end{equation}
that are source-independent, time-independent and common to every transition in the dataset.  ``Physical'' refers to complete positivity together with preservation of the physical trace $\tau$ of the represented record algebra.  The map is defined on the whole algebra; computational support reductions are permitted only when accompanied by a CPTP extension/retraction argument establishing minimax equivalence to this whole-algebra problem.

For each transition define
\begin{equation}
e_j(\Phi)
=\frac12\left\|
\Lambda_j^{\mathrm{next}}
-\Phi\circ\Lambda_j^{\mathrm{cur}}
\right\|_\diamond .
\label{eq:error}
\end{equation}
The diamond norm retains the complete source-channel scope and the explicit factor $1/2$ is part of the frozen error convention.

For a finite transition dataset $D_K$ of cardinality $K$, define the minimax incompatibility
\begin{equation}
\eps_*^{(K)}
=\inf_{\Phi\in\C}\max_{j\in D_K}e_j(\Phi),
\label{eq:minimax}
\end{equation}
and, at tolerance $\eps\ge0$, the feasible set
\begin{equation}
\F_K(\eps)
=\left\{\Phi\in\C:
\max_{j\in D_K}e_j(\Phi)\le\eps
\right\}.
\label{eq:feasible-set}
\end{equation}
The compatibility question at tolerance $\eps$ is simply whether $\F_K(\eps)$ is nonempty.  It is logically distinct from whether the finite data identify a unique or predictively stable representative.

\subsection{Accumulating constraints}

The useful ordering variable is not a solver objective from one run but the minimax value of the full physical class under a genuinely nested data sequence.

\begin{lemma}[Nested feasible sets]
\label{lem:nesting}
If $D_K\subseteq D_{K+r}$ for $r>0$, with the same admissible class $\C$, error definition and tolerance, then
\begin{equation}
\F_{K+r}(\eps)\subseteq\F_K(\eps).
\label{eq:set-nesting}
\end{equation}
Consequently,
\begin{equation}
\eps_*^{(K)}\le\eps_*^{(K+r)}.
\label{eq:minimax-monotonicity}
\end{equation}
\end{lemma}

\begin{proof}
If $\Phi\in\F_{K+r}(\eps)$, then $e_j(\Phi)\le\eps$ for every $j\in D_{K+r}$.  Since $D_K\subseteq D_{K+r}$, the same inequalities hold for every $j\in D_K$, hence $\Phi\in\F_K(\eps)$.  For the minimax values, for every fixed $\Phi$,
\[
\max_{j\in D_K}e_j(\Phi)
\le
\max_{j\in D_{K+r}}e_j(\Phi).
\]
Taking the infimum over the same class $\C$ preserves the inequality.
\end{proof}

The lemma is elementary, but it prevents a recurrent interpretive error.  A particular witness may leave the feasible set when new constraints are added without the enlarged feasible set becoming empty.  Specifically,
\begin{equation}
\Phi_K\notin\F_{K+r}(\eps)
\qquad\not\Rightarrow\qquad
\F_{K+r}(\eps)=\varnothing.
\label{eq:witness-not-class}
\end{equation}
This distinction is the central logical point of the present example.

\subsection{Empty-feasible-set semantics}

A nonempty feasible set describes models compatible with the sampled data at the declared tolerance.  An empty feasible set has the opposite meaning: the entire model class is incompatible with those constraints at that tolerance.  Accordingly, geometric quantities such as $\diam\F_K(\eps)$ are left undefined (or explicitly marked inapplicable) when $\F_K(\eps)=\varnothing$.  We do not assign the empty set diameter zero and interpret this as perfect identification.

\subsection{Three questions}

We distinguish three levels throughout:

\paragraph{Class compatibility.}
Does there exist at least one $\Phi\in\C$ satisfying all observed transition constraints?  This is equivalent to $\F_K(\eps)\neq\varnothing$.

\paragraph{Witness identification or selection.}
How tightly do the finite observations constrain the map, and by what rule is a representative $\Phi_K$ selected?  A feasible point is not necessarily unique, minimax-optimal, statistically optimal, or stable outside the constrained domain.

\paragraph{Prospective prediction.}
Does a previously frozen $\Phi_K$ accurately reproduce genuinely unseen transitions without refitting?  This is an out-of-sample test of that fixed witness, not an all-map existence test unless an additional universal certificate is supplied.

No implication among these questions should be assumed beyond what is proved from the data and class definitions.
